\documentclass[11pt, a4paper]{article}

%% >>> preamble
\usepackage[T1]{fontenc}
\usepackage{microtype}
\usepackage[a4paper, margin=2.5cm]{geometry}
\usepackage{amsmath, amssymb}
\usepackage{graphicx}
\usepackage{booktabs}
\usepackage{xcolor}
\usepackage[colorlinks=true, allcolors=blue!60!black]{hyperref}
\usepackage[capitalise, noabbrev]{cleveref}
%% <<< preamble

\title{A Short Article Title}
\author{Your Name\\[2pt]
  \small Your Institution\\
  \small\texttt{you@example.com}}
\date{\today}

\begin{document}
%% >>> snippet
\maketitle

\begin{abstract}
  One paragraph: the problem, what you did, what you found, and why it
  matters. Keep it to around 150 words.
\end{abstract}

\section{Introduction}\label{sec:intro}
Motivate the problem and place it in context~\cite{knuth1984}.
Our contributions are:
\begin{itemize}
  \item the first contribution;
  \item the second contribution.
\end{itemize}

\section{Method}\label{sec:method}
Describe the approach. Equations are numbered and referenced with
\cref{eq:model}:
\begin{equation}\label{eq:model}
  y = f(x; \theta) + \varepsilon, \qquad \varepsilon \sim \mathcal{N}(0, \sigma^2).
\end{equation}

\section{Results}\label{sec:results}
\Cref{tab:results} summarises the results.
\begin{table}[htbp]
  \centering
  \caption{Main results.}
  \label{tab:results}
  \begin{tabular}{@{}lcc@{}}
    \toprule
    Method   & Accuracy (\%) & Time (s) \\
    \midrule
    Baseline & 81.2          & 12.0 \\
    Ours     & \textbf{88.3} & \textbf{9.4} \\
    \bottomrule
  \end{tabular}
\end{table}

\section{Conclusion}\label{sec:conclusion}
Summarise what was shown, then point to future work.

\begin{thebibliography}{9}
  \bibitem{knuth1984} D.~E. Knuth. \emph{The \TeX book}. Addison-Wesley, 1984.
\end{thebibliography}
%% <<< snippet
\end{document}
